\section{The endomorphism \texorpdfstring{$\Omega$}{Omega} of a complex of representations}\label{sec:setting}

\subsection{Complexes of representations and metrics}
Let $Q$ be a finite quiver with set of vertices $Q_0$ and set of arrows
$Q_1$; we write $a:i\to j$ for an arrow from $i$ to $j$. Loops and oriented
cycles are allowed. We write $K_m$ for the generalized Kronecker quiver with
vertices $0,1$ and $m\geq1$ arrows $a_1,\dots,a_m$ from $0$ to $1$. Fix
\emph{central charges} $z_i$, $i\in Q_0$, in the open upper half-plane, and
put $\alpha_i=\arg z_i\in(0,\pi)$; they define the central charge of a stability
condition on the bounded derived category of finite-dimensional
representations of $Q$ whose heart is the category of representations
\cite[Propositions~2.4 and~5.3]{bri07}. Phases in this paper are arguments, in radians; in the
normalization of \cite{bri07} they are divided by $\pi$, and a shift by one
degree changes them by $\pi$. A \emph{complex of representations}
$\mathsf E$ consists of finite-dimensional complex vector spaces $E_i^n$
($i\in Q_0$, $n\in\mathbb Z$, almost all zero), differentials
$d_i:E_i^n\to E_i^{n+1}$ with $d_i^2=0$, and for each arrow $a:i\to j$ a chain
map $\rho(a):E_i^\bullet\to E_j^\bullet$ of degree zero. Its central charge is
\[
 Z(\mathsf E)=\sum_{i,n}(-1)^nz_i\dim E_i^n .
\]
A \emph{metric} $h$ on $\mathsf E$ is a Hermitian metric $h_i^n$ on each
$E_i^n$; a complex with a metric is called metrized. The adjoint of a linear
map for the metrics is written ${}^\dagger$, inner products are
conjugate-linear in the first argument, and
$\norm T_{\HS}=(\tr T^\dagger T)^{1/2}$ is the Hilbert--Schmidt norm. For an
endomorphism $T$ of a Hermitian space, $\frac12(T+T^\dagger)$ is its
Hermitian part and $\Ima T=\frac1{2\ii}(T-T^\dagger)$ its imaginary part. We
write $E_i=\bigoplus_nE_i^n$
and $E=\bigoplus_iE_i$. The \emph{number operator} $\Nop$ is the endomorphism
of $E$ equal to $n$ on $E_i^n$, and $\gr=(-1)^{\Nop}$ is the grading. A
tangent vector to the space of metrics at $h$ is an endomorphism
$\xi=h^{-1}\delta h$ of $E$ that preserves each $E_i^n$ and is Hermitian for
$h$. In the algebra $\mathcal A(Q)$ below we write $\sqrt{-1}$ for the
imaginary unit, as in \cite{hkkp26}; elsewhere it is $\ii$, and the italic
$i$ is a vertex.

\subsection{The definition of Haiden, Katzarkov, Kontsevich and Pandit}
The graded $*$-algebra $\mathcal A(Q)$ of \cite[Definition~5.13]{hkkp26} has
the basis $e_i,f_i$ ($i\in Q_0$) in degree $0$, $a\in Q_1$ in degree $1$ and
$a^*$ in degree $-1$, and the nonzero products of basis elements are
\begin{gather*}
 e_i^2=e_i,\qquad e_if_i=f_ie_i=f_i,\qquad e_ja=ae_i=a,\qquad e_ia^*=a^*e_j=a^*,\\
 aa^*=\sqrt{-1}\,f_j,\qquad a^*a=-\sqrt{-1}\,f_i
\end{gather*}
for $a:i\to j$. In the graded tensor product
$\mathcal A(Q)\otimes\End(E)$ the product is
$(x\otimes T)(y\otimes S)=(-1)^{|T||y|}xy\otimes TS$. The involution of
$\mathcal A(Q)$ fixes the $e_i$ and $f_i$ and exchanges $a$ and $a^*$; on
$\End(E)$ it is $T^*=T^\dagger$ for even $T$ and $T^*=\sqrt{-1}\,T^\dagger$
for odd $T$, the adjoints being those of the metric. Let
$\mathcal B\subset\mathcal A(Q)\otimes\End(E)$ be the subalgebra
\[
 \mathcal B=\bigoplus_i\bigl(e_i\otimes\End(E_i)\oplus f_i\otimes\End(E_i)\bigr)
 \oplus\bigoplus_{a:i\to j}\bigl(a\otimes\Hom(E_i,E_j)\oplus a^*\otimes\Hom(E_j,E_i)\bigr),
\]
with unit $\sum_ie_i\otimes\id_{E_i}$. Put
\[
 A=\sum_ie_i\otimes d_i+\sum_aa\otimes\rho(a),\qquad
 F=(A-A^*)^2+\sqrt{-1}\sum_if_i\otimes z_i\id_{E_i} .
\]
Here $A^*=\sqrt{-1}\sum_ie_i\otimes d_i^\dagger+\sum_aa^*\otimes\rho(a)^\dagger$,
and this is the only use of the involution. These are elements of
$\mathcal B$. Following \cite[\S5.4]{hkkp26} we put
\begin{equation}\label{eq:Omegadef}
 \Omega:=\gr\,\bigl[\exp(-\sqrt{-1}\,F)\bigr]_f\in\End(E),
\end{equation}
where the exponential is taken in $\mathcal B$ and, for
$X=\sum_i(e_i\otimes X_{e_i}+f_i\otimes X_{f_i})+\dots\in\mathcal B$, the
endomorphism $[X]_f=\bigoplus_iX_{f_i}$ is the component along the $f_i$.
With
\begin{equation}\label{eq:DeltaMGamma}
 \Delta_i=d_id_i^\dagger+d_i^\dagger d_i,\qquad
 M_a=d_j^\dagger\rho(a)-\rho(a)d_i^\dagger,\qquad
 \Gamma_i=z_i+\sum_{a:i\to\cdot}\rho(a)^\dagger\rho(a)-\sum_{a:\cdot\to i}\rho(a)\rho(a)^\dagger ,
\end{equation}
where $M_a:E_i^\bullet\to E_j^{\bullet-1}$ for $a:i\to j$, one has
\cite[display after Remark~5.15]{hkkp26}
\begin{equation}\label{eq:F}
 -\sqrt{-1}\,F=-\sum_ie_i\otimes\Delta_i+\sum_aa\otimes M_a
 +\sqrt{-1}\sum_aa^*\otimes M_a^\dagger+\sum_if_i\otimes\Gamma_i .
\end{equation}
For a representation, that is, a complex in degree zero, $\Omega=\Gamma$:
at each vertex $\Omega$ is $z_i$ shifted by the moment map of the action of
the unitary group of $E_i$.

In \cite[Remark~5.14]{hkkp26} the algebra $\mathcal A(Q)$ is compared with
the differential forms on a Riemann surface: the $e_i$ correspond to
functions, $a$ and $a^*$ to forms of type $(1,0)$ and $(0,1)$, and the $f_i$
to two-forms. In this dictionary $A-A^*$ plays the role of a superconnection
in the sense of Quillen \cite{qui85} on the graded space $E$, $F$ that of
its curvature shifted by the central term, $[\,\cdot\,]_f$ that of the
component of top degree, and the functional $\tau$ of
Section~\ref{sec:primitive} that of the integral of the supertrace of this
component. Section~\ref{sec:bundles} carries the dictionary out for bundles.

\subsection{An integral formula}
For Hermitian $\Delta\geq0$ and an operator $G$ put
$K_\Delta(G)=\int_0^1e^{-(1-x)\Delta}Ge^{-x\Delta}dx$, and for operators
$X:E'\to E''$, $Y:E''\to E'$ and Hermitian $\Delta',\Delta''\geq0$ on $E'$,
$E''$,
\[
 \mathcal J_{\Delta'',\Delta'}(X,Y)
 =\int_{\substack{x,y\geq0\\x+y\leq1}}e^{-x\Delta''}X\,e^{-(1-x-y)\Delta'}\,Y\,e^{-y\Delta''}\,dx\,dy .
\]
In eigenbases of $\Delta$, $\Delta''$ and $\Delta'$, with eigenvalues
$\ell_i$, $\ell''_i$ and $\ell'_\mu$, these are given by divided differences
of the function $\mathsf f(x)=e^{-x}$, written $\mathsf f[\,\cdot\,]$, by the
Hermite--Genocchi formula:
\[
 K_\Delta(G)_{ik}=-\mathsf f[\ell_i,\ell_k]\,G_{ik},\qquad
 \mathcal J_{\Delta'',\Delta'}(X,Y)_{ik}=\sum_\mu\mathsf f[\ell''_i,\ell'_\mu,\ell''_k]\,X_{i\mu}Y_{\mu k},
\]
with $-\mathsf f[a,b]>0$ and $\mathsf f[a,b,c]>0$, by the mean value form of
divided differences. The expansion
$e^{X+Y}=e^X+\int_0^1e^{(1-x)X}Ye^{xX}dx+\dots$ is called the Duhamel
expansion below.

\begin{proposition}\label{prop:coefficient}
The endomorphism $\Omega$ preserves each $E_i^n$, and for every vertex $i$
\[
 \gr\,\Omega\big|_{E_i}=K_{\Delta_i}(\Gamma_i)
 +\sum_{a:k\to i}\mathcal J_{\Delta_i,\Delta_k}(M_a,M_a^\dagger)
 -\sum_{a:i\to j}\mathcal J_{\Delta_i,\Delta_j}(M_a^\dagger,M_a).
\]
\end{proposition}

\begin{proof}
Write the right side of \eqref{eq:F} as $X_e+X_a+X_{a^*}+X_f$. In
$\mathcal A(Q)$ every product of three elements of $\{a,a^*,f_i\}$ vanishes,
and so do the products $f_if_k$, $f_ia$, $af_i$, $f_ia^*$, $a^*f_i$, $aa'$,
$a^*a'^*$, and $aa'^*$, $a^*a'$ for $a\neq a'$. So in the Duhamel expansion of
$\exp(X_e+Y)$, $Y=X_a+X_{a^*}+X_f$, the component along the $f_i$ receives
only the terms with one factor $X_f$, or with one factor $X_a$ and one
factor $X_{a^*}$ of the same arrow, in either order. The element $X_e$ is
even and the $e_i$ act as units, so
$e^{xX_e}=\sum_ie_i\otimes e^{-x\Delta_i}$ and no signs arise from it. By the
sign rule, for $a:i\to j$,
\begin{align*}
 (a\otimes M_a)(\sqrt{-1}\,a^*\otimes M_a^\dagger)
 &=\sqrt{-1}\,(-1)^{|M_a||a^*|}aa^*\otimes M_aM_a^\dagger=f_j\otimes M_aM_a^\dagger,\\
 (\sqrt{-1}\,a^*\otimes M_a^\dagger)(a\otimes M_a)
 &=\sqrt{-1}\,(-1)^{|M_a^\dagger||a|}a^*a\otimes M_a^\dagger M_a=-f_i\otimes M_a^\dagger M_a .
\end{align*}
The second-order term of the expansion is
$\int e^{xX_e}Ye^{(1-x-y)X_e}Ye^{yX_e}\,dx\,dy$ over the simplex; inserting the
heat factors gives the three terms. Each term preserves the degree, because
$\Delta_i,\Gamma_i$ have degree $0$, $M_a$ has degree $-1$ and $M_a^\dagger$
degree $+1$.
\end{proof}

The formula is the expansion to second order of the display for
$-\sqrt{-1}F$ in \cite[\S5.4]{hkkp26}; no higher terms occur.

\begin{proposition}\label{prop:properties}
For every metrized complex:
\begin{enumerate}
\item[(a)] All terms of Proposition~\ref{prop:coefficient} other than the
 summand $K_{\Delta_i}(z_i)=z_ie^{-\Delta_i}$ of $K_{\Delta_i}(\Gamma_i)$ are
 Hermitian. Hence, on each $E_i^n$,
 $\Ima\bigl((-1)^n\Omega\bigr)=(\Ima z_i)\,e^{-\Delta_i}>0$.
\item[(b)] $\tr\Omega=Z(\mathsf E)$ \cite[Lemma~5.17]{hkkp26}.
\item[(c)] On each $E_i^n$ there is a unique Hermitian $\Phi_i^n$ with
 spectrum in $(-\pi n,\pi(1-n))$ such that
 $\Omega|_{E_i^n}\,e^{-\ii\Phi_i^n}$ is positive definite.
\item[(d)] The flow $h^{-1}\dot h=-\Phi$, $\Phi=\bigoplus\Phi_i^n$, has a
 unique solution for all $t\geq0$ from every initial metric.
\end{enumerate}
\end{proposition}

\begin{proof}
(a) $K_\Delta(G)^\dagger=K_\Delta(G^\dagger)$ by the substitution
$x\mapsto1-x$, and $\mathcal J(X,X^\dagger)^\dagger=\mathcal J(X,X^\dagger)$ by
exchanging $x$ and $y$. (b) This is \cite[Lemma~5.17]{hkkp26}; the argument
is as follows. With $\tau(X)=\tr(\gr\,[X]_f)$ one has
$\tr\Omega=\tau(e^{-\sqrt{-1}F})$. Replace all structure maps by their
multiples by $u\in[0,1]$; then $F$ becomes $-u^2[A,A^*]+F_z$ with
$F_z=\sqrt{-1}\sum_if_i\otimes z_i$, the element $A$ commutes with it, and the
derivative of $\tau(e^{-\sqrt{-1}F})$ in $u$ is
$2\sqrt{-1}\,u\,\tau\bigl([A,A^*e^{-\sqrt{-1}F}]\bigr)=0$, because $\tau$
vanishes on supercommutators of elements of $\mathcal B$
(Lemma~\ref{lem:tau} below). At $u=0$,
$e^{-\sqrt{-1}F_z}=\sum_i(e_i+f_i\otimes z_i)$ and
$\tau=\sum_{i,n}(-1)^nz_i\dim E_i^n$. (c) Let $T$ be a matrix with $\Ima T$ positive
definite. Then $T$ is invertible, since $Tv=0$ gives
$\inner v{(\Ima T)v}=\Ima\inner v{Tv}=0$. Let $T=PU$ be its polar
decomposition, with $P=(TT^\dagger)^{1/2}>0$ and $U$ unitary. If
$Uv=\lambda v$, then $\inner v{(\Ima T)v}=(\Ima\lambda)\inner v{Pv}$, so
$\Ima\lambda>0$. Hence $U=e^{\ii\Psi}$ for a unique Hermitian $\Psi$ with
spectrum in $(0,\pi)$, and $Te^{-\ii\Psi}=P>0$; conversely such a $\Psi$
gives a polar decomposition, which is unique. Apply this to
$(-1)^n\Omega|_{E_i^n}$ and put $\Phi_i^n=\Psi-\pi n$. (d) The vector field is real-analytic on the space of
metrics, and $\norm{h^{-1}\dot h}\leq\pi(1+\max|n|)$, so solutions cannot
leave compact sets in finite time.
\end{proof}

The computation in (a) is that of \cite[proof of Lemma~5.16]{hkkp26}. That
lemma also states that $\Omega$ is normal. By (a), a block $\Omega|_{E_i^n}$
is normal if and only if its Hermitian part commutes with $e^{-\Delta_i}$,
and this does not hold for all complexes:

\begin{example}[$\Omega$ need not be normal]\label{ex:nonnormal}
Let $Q=K_1$, with charges $z_0=-1+\ii$, $z_1=1+\ii$, and consider the complex
with the spaces $V=\CC^2$ in degree $0$ and $U_0=\CC^2$ in degree $1$ at the
vertex $0$, the spaces $U_1=\CC$ in degree $-1$ and $W=\CC$ in degree $0$ at
the vertex $1$, the standard metrics, the differentials
$d_0=\diag(1,4):V\to U_0$ and $d_1=3:U_1\to W$, and the arrow
$\rho(a_1)=(\rho_1,\rho_2)=(4,100):V\to W$. On $V$ one has
$d_0^\dagger d_0=\diag(\ell_1,\ell_2)$ with $(\ell_1,\ell_2)=(1,16)$ and
$\Gamma_0=z_0+\rho(a_1)^\dagger\rho(a_1)$; the operator $M_{a_1}$ of
\eqref{eq:DeltaMGamma} restricts to $3\rho(a_1):V\to U_1$, and $\Delta_1=9$
on $U_1$. By Proposition~\ref{prop:coefficient} (see
Proposition~\ref{prop:twosidedcoefficient} for the general formula), the
block of $\Omega$ on $V$ is
\[
 (\Omega_V)_{ik}=z_0e^{-\ell_i}\delta_{ik}
 +\rho_i\rho_k\bigl(-\mathsf f[\ell_i,\ell_k]-9\,\mathsf f[\ell_i,9,\ell_k]\bigr),
\]
where the factor $9$ is $|d_1|^2$ and the point $9$ is the eigenvalue of
$d_1^\dagger d_1$. Its imaginary part $\diag(e^{-1},e^{-16})$ has distinct
eigenvalues, and its Hermitian part has the off-diagonal entry
$400\,(-\mathsf f[1,16]-9\,\mathsf f[1,9,16])=-1.2183\ldots\neq0$; so the two
do not commute. This complex is used again in
Theorem~\ref{thm:twosidedfailure}. For two-term complexes over $K_m$ see
\cite[Example~2.3]{wue26a}.
\end{example}

\subsection{Terminology}

\begin{definition}\label{def:phase}
Let $\mathsf E$ be a metrized complex.
\begin{enumerate}
\item[(i)] The \emph{blocks} of $\Omega$ are the endomorphisms
 $\Omega_i^n=\Omega|_{E_i^n}$. The \emph{phase} is
 $\Phi=\bigoplus\Phi_i^n$ with $\Phi_i^n$ as in
 Proposition~\ref{prop:properties}(c), and
 $|\Omega_i^n|:=\Omega_i^ne^{-\ii\Phi_i^n}=(\Omega_i^n\Omega_i^{n\dagger})^{1/2}$.
 The word phase is also used, as usual, for the argument of a central
 charge.
\item[(ii)] The \emph{phase flow} is the flow $h^{-1}\dot h=-\Phi$ of
 Proposition~\ref{prop:properties}(d). The \emph{mass} is
 $\mass(h)=\sum_{i,n}\tr|\Omega_i^n|$, and a metric is \emph{harmonic} if
 $\Phi$ is a scalar.
\item[(iii)] Let $i$ be a vertex. The complex of vector spaces $E_i^\bullet$
 splits orthogonally as $\ker\Delta_i\oplus(\ker\Delta_i)^\perp$; the
 differential vanishes on the first summand, which maps isomorphically onto
 the cohomology, and the second summand is an acyclic subcomplex. We call
 $(\ker\Delta_i)^\perp$ the \emph{acyclic part of $\mathsf E$ at $i$} and the
 rank of $d_i$ its \emph{rank}. The metrics \emph{induced on the cohomology}
 are the restrictions of the metrics to the spaces $\ker\Delta_i$.
\item[(iv)] A \emph{contractible summand} is a direct summand of
 $\mathsf E$, as a complex of representations, that is contractible. Those
 used below are complexes $C^n\to C^{n+1}$ at one vertex, with invertible
 differential, on which no arrow acts.
\item[(v)] Suppose that the cohomology of $\mathsf E$ is a representation in
 degree zero, with $Z(\mathsf E)=|Z|e^{\ii\varphi}$, $\varphi\in(0,\pi)$.
 Along a trajectory of the phase flow, the acyclic part \emph{decays} if the
 least nonzero eigenvalue of each $\Delta_i$ tends to infinity and $\Omega$
 minus the endomorphism of the cohomology representation with its induced
 metrics, extended by zero to the acyclic parts, tends to zero in the trace
 norm $\tr(T^\dagger T)^{1/2}$ for the metrics $h(t)$ (this is called
 cancellation of the acyclic part
 in \cite{wue26b} and decay of the exact part in \cite{wue26c}); and the
 flow
 \emph{converges on the cohomology} if the metrics induced on the
 cohomology, multiplied by $e^{\varphi t}$, converge to a harmonic metric of
 the cohomology representation.
\end{enumerate}
\end{definition}

\begin{remark}\label{rem:phase}
(1) The phase of Definition~\ref{def:phase} satisfies the positivity
condition of \cite[\S4.2]{hkkp26} recalled in the Introduction, with the
branches of \cite[(5.15)]{hkkp26}, and the phase flow is the minimizing flow
of \cite[\S4.1]{hkkp26} for this phase; for normal $\Omega$ the phase is the
argument in the sense of the functional calculus used in
\cite[\S5.4]{hkkp26}.

(2) Since the intervals $(-\pi n,\pi(1-n))$ are pairwise disjoint, a harmonic
metric can exist only on a complex concentrated in one degree $n$, that is,
on a shifted representation; then $\Phi$ is the scalar $\arg Z(\mathsf E)$
taken in $(-\pi n,\pi(1-n))$, the metric is a zero of the moment map shifted
by the stability parameter (Proposition~\ref{prop:kempfness}), and one
exists if and only if the representation is polystable
\cite[Theorem~5.1]{hkkp26}, \cite[Proposition~6.5]{kin94}. For a complex with
terms in several degrees, `harmonic' is used below only for the metrics
induced on the cohomology.

(3) The acyclic part at a vertex is in general not preserved by the arrows,
and an acyclic direct summand need not be contractible
(Theorem~\ref{thm:twosidedfailure}).
\end{remark}

Changing the metric at fixed structure maps is the same as moving the
structure maps in their orbit under $\prod\mathrm{GL}(E_i^n)$ at a fixed
metric. Precisely, if $h=\overline g^{\mathsf T}g$ in fixed bases, the matrix
of a structure map
$T\in\{d_i,\rho(a)\}$ in the orthonormal frames given by $g$ is
$g_{\mathrm{target}}Tg_{\mathrm{source}}^{-1}$; along the phase flow one can
choose $g$ with $\dot gg^{-1}=-\frac12\Phi$, where $\Phi$ now denotes the
Hermitian matrix of the phase in these frames, and then
\begin{equation}\label{eq:frameflow}
 \dot T=\tfrac12\bigl(T\,\Phi_{\mathrm{source}}-\Phi_{\mathrm{target}}\,T\bigr)
\end{equation}
for the matrices in orthonormal frames, where $\Phi_{\mathrm{source}}$,
$\Phi_{\mathrm{target}}$ are the phases on the source and the target of $T$.
The endomorphism $\Omega$, the phases and the mass depend only on the
matrices of the structure maps in orthonormal frames, up to unitary changes
of the frames.

For $Q=K_m$ the charges are called \emph{ordered} if
$\alpha_0\geq\alpha_1$ and \emph{strictly ordered} if $\alpha_0>\alpha_1$;
that is, the simple representation at the source vertex has phase at least,
respectively larger than, that of the simple representation at the target
vertex. We
write a complex of representations of $K_m$ as a row of its terms, each term
being the representation $(E_0^n\to E_1^n)$ with the space at the vertex $0$
on the left; the maps between consecutive terms are the differentials
$(d_0,d_1)$.

\subsection{Two matrix inequalities}
We recall two matrix inequalities of \cite{wue26b} and strengthen the first. For Hermitian matrices $L>0$
and $M$ put
\[
 J_L(M)=\int_{\substack{x,y\geq0\\x+y\leq1}}e^{-xL}Me^{-yL}\,dx\,dy ,
\]
so that $J_L(XY)=\mathcal J_{L,0}(X,Y)$; in an eigenbasis of $L$ with
eigenvalues $\ell_i$, $J_L(M)_{ik}=\mathsf f[\ell_i,0,\ell_k]\,M_{ik}$. We
write $M_\pm\geq0$ for the positive and negative parts of a Hermitian matrix
$M=M_+-M_-$.

\begin{theorem}[{\cite[Theorems~3.1 and~4.1]{wue26b}}]\label{thm:imported}
The following hold.
\begin{enumerate}
\item[(a)] For Hermitian matrices $L>0$ and $M\geq0$,
 \begin{equation}\label{eq:kernel}
 \tr J_L(M)_-\leq\tfrac25\tr J_L(M).
 \end{equation}
\item[(b)] Let $T(s)$ be a smooth family of matrices with constant positive
 definite imaginary part, and $\Psi(s)$ the Hermitian matrices with spectrum
 in $(0,\pi)$ and $T e^{-\ii\Psi}>0$. Then $\tr(\dot T\,\dot\Psi)<0$ wherever
 $\dot T\neq0$.
\end{enumerate}
\end{theorem}

For $1\times1$ matrices, (a) is empty because $J_L(M)\geq0$, and (b) says
that $\arg(x+\ii y)$ decreases in $x$ for fixed $y>0$. In \cite{wue26b}, (a)
is reduced to an inequality in one real variable, proved with the help of
interval arithmetic, and (b) is proved by expressing
$\tr(\dot T\,\dot\Psi)$ through a Sylvester equation as a quadratic form whose
kernels are positive semidefinite by an integral representation
\cite[Section~4]{wue26b}.

We shall use the stronger estimate
$\tr J_L(M)_-\leq\frac13\tr J_L(M)$, which admits a direct proof.
In an eigenbasis of $L$, write $J_L(M)_{ik}=M_{ik}j(\ell_i,\ell_k)$,
where $j(a,b)$ is the integral of $e^{-ax-by}$ over $x,y\geq0$, $x+y\leq1$.
For $s\in I=[0,1/2]$ put $F_a(s)=\int_0^s e^{-ax}\,dx$,
$A_a=F_a(1/2)$ and $V_a(s)=e^{-a(1-s)}$. Splitting the triangle at
$x=1/2$ and $y=1/2$ gives
$j(a,b)=A_aA_b+\int_I(F_aV_b+V_aF_b)\,ds$. Thus both $Q$ and $j+Q$ are
positive semidefinite kernels, where
\[
 Q(a,b)=\tfrac23\int_I(F_a-\tfrac34V_a)(F_b-\tfrac34V_b)\,ds .
\]
Indeed, $j+Q$ is $A_aA_b$ plus the same integral with both minus signs
replaced by plus signs. On the diagonal suppress the subscript $a$. Since
$sV\leq F\leq A$ and $V^2$ is increasing,
$\int_I FV\geq\int_I sV^2\geq\tfrac14\int_I V^2$. Expanding the square gives
\[
 Q(a,a)\leq\tfrac13 A^2+\tfrac12\int_I FV\,ds
 \leq\tfrac13 j(a,a).
\]
By the Schur product theorem, the matrices with entries
$M_{ik}(j+Q)(\ell_i,\ell_k)$ and $M_{ik}Q(\ell_i,\ell_k)$ are positive
semidefinite. Taking the trace on the negative spectral subspace of their
difference bounds $\tr J_L(M)_-$ by $\sum_i M_{ii}Q(\ell_i,\ell_i)$,
which is at most $\frac13\tr J_L(M)$.
This estimate enters Theorem~\ref{thm:onevertex}(b),(c) in rank at least
two and hence the one-sided cases of Corollary~\ref{cor:classification}.
Proposition~\ref{prop:parabolic} uses Theorem~\ref{thm:imported}(b).
